\section{Thiết kế \& Triển khai Smart Contract (Solidity --- Backend phi tập trung)}
\label{sec:ch5}

Smart contract \code{TourReservation} (tệp \code{TourReservation.sol}, giấy phép MIT, \code{pragma solidity \^{}0.8.20})
đóng vai trò ``backend phi tập trung'': nó vừa lưu dữ liệu, vừa thực thi các quy tắc nghiệp vụ, vừa giữ tiền của khách.
Toàn văn mã nguồn có ở \ref{sec:appendix}.

\subsection{Biến trạng thái và cấu trúc dữ liệu}

\lstinputlisting[style=SolStyle, firstline=4, lastline=24,
  caption={Biến trạng thái, mapping, sự kiện, modifier và hàm khởi tạo}]{Code/TourReservation.sol}

\begin{table}[H]
\centering
\caption{Biến trạng thái của hợp đồng}
\begin{tabularx}{\linewidth}{|l|l|X|}
\hline
\rowcolor{navyblue}\textcolor{white}{\textbf{Biến}} & \textcolor{white}{\textbf{Kiểu}} & \textcolor{white}{\textbf{Ý nghĩa}} \\
\hline
\code{operator} & \code{address public} & Ví điều hành, gán bằng \code{msg.sender} của người triển khai trong \code{constructor}. \\
\hline
\rowcolor{rowgray}
\code{totalSeats} & \code{uint256 public} & Tổng số chỗ của tour (ví dụ 25). \\
\hline
\code{availableSeats} & \code{uint256 public} & Số chỗ còn trống; giảm dần khi có người đặt. \\
\hline
\rowcolor{rowgray}
\code{ticketPrice} & \code{uint256 public} & Giá một vé tính bằng \emph{wei}; 0,08 ETH $=80\,000\,000\,000\,000\,000$ wei $= 8\times10^{16}$. \\
\hline
\code{isTourActive} & \code{bool public} & \code{true} sau khi điều hành gọi \code{createTour()}; chặn đặt chỗ khi chưa có tour. \\
\hline
\rowcolor{rowgray}
\code{bookings} & \code{mapping(address => uint256) private} & \textbf{Số vé đã đặt theo từng ví}. Khóa là địa chỉ ví, giá trị là số chỗ. \\
\hline
\end{tabularx}
\end{table}

\textbf{Mapping lưu vé theo ví.} \code{mapping(address => uint256) private bookings} là bảng băm trên chuỗi: với mỗi địa chỉ,
giá trị mặc định là 0 và tăng lên khi ví đó đặt chỗ. Từ khóa \code{private} nghĩa là các hợp đồng khác không đọc trực tiếp được
biến này; việc xem số vé thực hiện qua hàm \code{getMyBookings()}. (Lưu ý: \code{private} chỉ giới hạn truy cập ở mức mã, dữ liệu trên
blockchain vẫn công khai về mặt kỹ thuật.)

\textbf{Sự kiện (event).} \code{TourCreated} được phát khi tạo tour; \code{SeatsBooked} (với \code{customer} được \code{indexed})
được phát mỗi lần đặt chỗ, giúp ứng dụng bên ngoài lọc lịch sử đặt chỗ theo ví mà không cần lưu thêm dữ liệu.

\textbf{Modifier \code{onlyOperator}.} Gắn vào hàm để kiểm tra \code{msg.sender == operator} trước khi chạy thân hàm
(\code{\_;}). Nếu sai, giao dịch bị hoàn tác (revert) với thông báo \emph{``Chi dieu hanh moi co quyen thuc hien''}.

\subsection{Hàm \code{createTour()}}

\lstinputlisting[style=SolStyle, firstline=26, lastline=37,
  caption={Hàm \code{createTour} --- chỉ điều hành, số chỗ và giá vé dương}]{Code/TourReservation.sol}

\begin{itemize}[leftmargin=1.5em]
    \item Có \code{onlyOperator}: chỉ ví điều hành gọi được (ràng buộc 1 của đề).
    \item \code{require(\_totalSeats > 0)} và \code{require(\_ticketPrice > 0)}: số chỗ và giá vé phải dương.
    \item Ghi \code{totalSeats}, \code{availableSeats = \_totalSeats}, \code{ticketPrice}, \code{isTourActive = true}
          và phát sự kiện \code{TourCreated}.
\end{itemize}

\subsection{Hàm \code{bookSeats()}}

\lstinputlisting[style=SolStyle, firstline=39, lastline=50,
  caption={Hàm \code{bookSeats} --- đặt chỗ và thanh toán}]{Code/TourReservation.sol}

Hàm có từ khóa \code{payable} nên nhận được ETH kèm giao dịch (\code{msg.value}). Bốn câu lệnh \code{require} theo đúng thứ tự kiểm tra:

\begin{table}[H]
\centering
\caption{Các điều kiện \code{require} của \code{bookSeats()}}
\begin{tabularx}{\linewidth}{|c|l|X|}
\hline
\rowcolor{navyblue}\textcolor{white}{\textbf{STT}} & \textcolor{white}{\textbf{Điều kiện}} & \textcolor{white}{\textbf{Thông báo khi vi phạm (revert)}} \\
\hline
1 & \code{isTourActive} & ``Hien chua co tour nao duoc mo'' \\
\hline
\rowcolor{rowgray}
2 & \code{\_seats > 0} & ``So cho dat phai lon hon 0'' \\
\hline
3 & \code{\_seats <= availableSeats} & ``Khong du cho trong'' \\
\hline
\rowcolor{rowgray}
4 & \code{msg.value == \_seats * ticketPrice} & ``Thanh toan khong dung phi'' \\
\hline
\end{tabularx}
\end{table}

Nếu cả bốn điều kiện đúng, hợp đồng cập nhật trạng thái: \code{availableSeats -= \_seats} (giảm chỗ trống) và
\code{bookings[msg.sender] += \_seats} (cộng dồn số vé của ví), rồi phát sự kiện \code{SeatsBooked}.
Vì \code{\_seats} kiểu \code{uint256} nên không thể âm; Solidity $\ge$ 0.8 tự động revert khi tràn số, nên phép nhân
\code{\_seats * ticketPrice} an toàn. Điều kiện 4 yêu cầu trả \emph{đúng} bằng, không cho trả thừa hay thiếu, nên hợp đồng
không cần tính tiền thừa.

\subsection{Hàm \code{getAvailableSeats()} và \code{getMyBookings()}}

\lstinputlisting[style=SolStyle, firstline=52, lastline=60,
  caption={Hai hàm đọc dữ liệu}]{Code/TourReservation.sol}

Cả hai là hàm \code{view} nên chạy miễn phí. \code{getMyBookings()} dùng \code{msg.sender} làm khóa tra \code{bookings};
mỗi ví chỉ xem được số vé \emph{của chính mình} thông qua hàm này.

\subsection{Quá trình Compile \& Deploy trên Remix IDE}

\subsubsection{Biên dịch (Compile)}

\begin{enumerate}[leftmargin=1.8em]
    \item Mở Remix IDE, tạo tệp \code{contracts/TourReservation.sol} và dán mã nguồn.
    \item Mở tab \emph{Solidity Compiler}, chọn trình biên dịch \textbf{0.8.34+commit.80d5c536} (phù hợp \code{\^{}0.8.20}), nhấn \emph{Compile}.
    \item Biên dịch thành công khi dấu tích xanh xuất hiện ở biểu tượng Compile.
\end{enumerate}

\shot{compile.png}{Remix IDE: biên dịch thành công \code{TourReservation.sol} và kết quả triển khai lên Sepolia (log xác minh mã nguồn trên Blockscout/Sourcify)}{0.95}

\subsubsection{Triển khai (Deploy)}

\begin{enumerate}[leftmargin=1.8em]
    \item Tab \emph{Deploy \& Run Transactions}: chọn \textbf{Environment = Browser Extension} (MetaMask, mạng Sepolia).
          Ví triển khai là \textbf{Account 3} --- ví này tự động trở thành \emph{operator}.
    \item Nhấn \emph{Deploy}. MetaMask hiện cửa sổ \emph{Deploy a contract} với phí mạng ước tính khoảng 0,0073 SepoliaETH; nhấn \emph{Confirm} để ký.
    \item Sau khi giao dịch được đưa vào khối, Remix hiển thị hợp đồng trong mục \emph{Deployed Contracts} cùng các hàm
          \code{bookSeats}, \code{createTour}, \dots\ để gọi thử.
\end{enumerate}

\shot{confirm_deloy.png}{MetaMask yêu cầu xác nhận triển khai hợp đồng (\emph{Deploy a contract}) trên Sepolia, phí mạng $\approx$ 0,0073 SepoliaETH}{0.95}

\begin{table}[H]
\centering
\caption{Thông tin triển khai}
\begin{tabularx}{\linewidth}{|l|X|}
\hline
\rowcolor{navyblue}\textcolor{white}{\textbf{Mục}} & \textcolor{white}{\textbf{Giá trị}} \\
\hline
Mạng & Ethereum Sepolia Testnet \\
\hline
\rowcolor{rowgray}
Trình biên dịch & Solidity \code{0.8.34+commit.80d5c536} (Remix 2.6.14) \\
\hline
Ví triển khai / điều hành (Account 3) & \code{0x04fd74c1feb8b77b5f8A6a3f1be546778D06BBa2} \\
\hline
\rowcolor{rowgray}
Địa chỉ hợp đồng dùng trong DApp & \code{0xAa5AfB26CFcc5Fdb6EedAeff80A6E17bD0D0902e} \\
\hline
Mã giao dịch (hash) lần triển khai trong ảnh chụp & \code{0xdb4d19e8fdfc830b3ceb\allowbreak 14269ff4d0f0fda9ef7d02\allowbreak 713423952dfaadb129a936} \\
\hline
\rowcolor{rowgray}
Khối (block) & 12139637 \\
\hline
Gas sử dụng khi deploy (\emph{transaction cost}) & 4\,858\,715 gas \\
\hline
\end{tabularx}
\end{table}

\shot{gas.png}{Chi tiết giao dịch triển khai trong Remix: trạng thái \emph{mined and execution completed}, mã giao dịch, số khối và chi phí 4\,858\,715 gas}{0.8}

Sau khi triển khai, Etherscan trên Sepolia ghi nhận hợp đồng cùng các giao dịch gọi hàm (Hình~\ref{fig:eth_smartcontract.png} ở \ref{sec:ch7}).
Trong quá trình phát triển, hợp đồng được triển khai nhiều lần (Remix liệt kê nhiều mục trong \emph{Deployed Contracts});
bản được dùng cho DApp là địa chỉ đã khai báo trong \code{CONTRACT\_ADDRESS}.

\begin{infobox}
\textbf{Sau khi deploy:} sao chép địa chỉ hợp đồng từ Remix, dán vào \code{CONTRACT\_ADDRESS} trong
\code{frontend/src/config/contract.js}. ABI chỉ cần cập nhật khi thay đổi chữ ký hàm của hợp đồng.
\end{infobox}
